%% ==========================================================================
%%  5  The Replica Instance
%%
%%  MAIN REPORT ONLY -- not included by working.tex.  Labels prefixed `main:'.
%%
%%  Carries the Approach 3 material that the main report needs (the replica
%%  construction, the transform, and coverage-is-the-entire-gap), and then the
%%  third obstruction, which is Approach 4: enforcing coverage by pricing.
%%
%%  Condensed from working/approach_3.tex (sec:replica, thm:replica,
%%  cor:coverage) and working/approach_4.tex (thm:w1, sec:holderblind).
%% ==========================================================================
\section{The Replica Instance}
\label{sec:replica}

The third of the techniques named in \Cref{sec:obstructions} turns the chore
problem into a \emph{goods} problem, to which the theorem of~\cite{BKNS22}
applies verbatim, and then enforces what is left over by pricing. The translation
is exact and is worth having in its own right; the pricing step is the third
obstruction, and is treated in \Cref{sec:main-pricing}.

% --------------------------------------------------------------------------
\subsection{The construction}
\label{sec:main-replica}

Fix an instance with cost functions $\cost_i$ on chores
$\items = \set{a_1,\dots,a_m}$. The \emph{replica instance} $\hat{I}$ carries
$n-1$ copies of a good $b_j$ for each chore $a_j$,
\[
  \hat{\items} \;:=\; \set{\, b_j^{(1)}, \dots, b_j^{(n-1)} \;:\; j \in \items \,},
\]
and for $S \subseteq \hat{\items}$, writing $\tau(S) \subseteq \items$ for the
set of \emph{types} occurring in $S$,
\[
  \hat{v}_i(S) \;:=\; \cost_i(\items) \;-\; \cost_i\big(\items \setminus \tau(S)\big).
\]

The reading is that \emph{handing agent $i$ a copy of $b_j$ relieves $i$ of chore
$a_j$}. Each chore ends with exactly one owner, so exactly $n-1$ agents must be
relieved of it --- hence $n-1$ copies --- and no agent may be relieved of the
same chore twice, so no bundle may hold two copies of one type. Call an
allocation of $\hat{\items}$ satisfying that last condition a \emph{coverage
allocation}. The condition is load-bearing rather than cosmetic: it is what makes
$\hat{v}_i$ well defined as a dual, a second copy of a type having marginal $0$
and never $+1$ again.

\begin{theorem}[Replica transform]
\label{main:replica}
Let $\cost_1,\dots,\cost_n$ be dichotomous. Then
\begin{enumerate}
  \item[(i)] each $\hat{v}_i$ is a dichotomous valuation on $\hat{\items}$ in the
    sense of~\cite{BKNS22};
  \item[(ii)] the map $\Phi$ sending an allocation $\alloc$ of $\items$ to the
    allocation in which bundle $i$ receives one copy of each type in
    $\items \setminus \bundle{i}$ is a bijection between partitions of $\items$
    and coverage allocations of $\hat{\items}$, up to relabelling copies of a
    type;
  \item[(iii)] corresponding allocations have identical envy graphs:
    $\hat{\edgew{}}_{B}(i,k) = \edgew{\alloc}(i,k)$ for all $i,k$.
\end{enumerate}
\end{theorem}

Since the envy graphs agree arc for arc, so do the feasible subsidy vectors, and
the $\set{0,1}$ target transfers in both directions without loss. The problem is
therefore not merely analogous to a dichotomous goods problem; on coverage
allocations it \emph{is} one.

\begin{corollary}[Coverage is the entire gap]
\label{main:coverage}
Conjecture~\ref{conj:main} holds if and only if every replica instance
$\hat{I}$ admits a coverage allocation $B$ with $\pathw{B}(i) \le 1$ for all $i$.
\end{corollary}

% --------------------------------------------------------------------------
\subsection{Pricing the coverage constraint}
\label{sec:main-pricing}

Applying~\cite{BKNS22} to $\hat{I}$ is legal by Theorem~\ref{main:replica}(i) and
returns some $B$ with $\hat{\subsidy} \in \set{0,1}^n$ and total at most $n-1$.
But it is free to place two copies of one type in one bundle, in which case that
chore is absent from at least two bundles, has at least two owners, and the image
under $\Phi^{-1}$ is not a partition. By Corollary~\ref{main:coverage} coverage
is the sole residual difficulty, so it suffices to make duplicates unattractive.

Two ways of doing so suggest themselves. One may \emph{stage} the allocation,
dealing out the $n-1$ copies of $a_1$ together, then those of $a_2$, and so on,
with each chore priced far above the last, so that an agent who has just taken a
copy of the current chore is too well supplied to be selected again while copies
of that chore remain; staging makes coverage automatic, since a stage placing its
$n-1$ copies on distinct agents yields a coverage allocation by construction.
Or one may simply charge for duplicates directly, adding a penalty proportional
to the number of repeated copies a bundle holds. Neither works, and both fail for
the same reason.

\begin{example}[Duplicates are invisible to any price]
\label{main:ex-pricing}
In the replica a second copy of a chore has marginal value exactly $0$: if $b$
has the type of a chore already represented in $S$ then
$\tau(S \cup \set{b}) = \tau(S)$, so $\hat{v}_i(S \cup \set{b}) = \hat{v}_i(S)$
for every agent $i$. The holder of a duplicate is therefore \emph{indifferent} to
it, not averse, which is the wrong sign: a selection rule is not repelled by the
duplicate, it merely gains nothing from avoiding it.

Nor can a price supply the missing aversion. A charge levied on duplicates bites
only through the arcs of the envy graph, and there it cancels: whatever an agent
is charged for the duplicates she holds lowers every arc \emph{out of} her by
exactly what the charge raises the arcs \emph{into} her from everyone else, so
the two contributions annihilate around every cycle, and by
Theorem~\ref{thm:hs-characterisation}(iii) no allocation is disqualified. A price
can therefore work only if the agent who \emph{holds} the duplicate does not feel
it while somebody else does --- and since prices are fixed before the allocation
is chosen, while the allocation decides who holds the duplicate, requiring
holder-blindness for every possible holder forces the price to zero.
\end{example}

Staging does not evade this. Its escalation separates \emph{chores}, whereas the
duplicate problem lives \emph{within} one chore: the two copies at issue are
copies of the same $a_j$, literally the same object, and no price assigned to
chores can tell them apart. A further obstacle stands behind that one, and is
fatal on its own --- making one chore ``much costlier'' than another requires
marginals outside $\set{0,1}$, which leaves the dichotomous class and forfeits
the very theorem the staging was meant to invoke.

% --------------------------------------------------------------------------
\subsection{What the three failures have in common}
\label{sec:main-expost}

Each of the three techniques fixes a rule before the quantity it must respond to
has been determined. Insertion fixes an owner before the consequences of that
grouping are visible, and cannot later re-split a bundle. A selection rule fixes
its criterion before knowing which allocation the criterion will pick out. A
price fixes a charge before knowing who will hold the duplicate being charged
for. In each case the rule is settled \emph{ex ante} against a quantity
determined \emph{ex post}.

What survives is the requirement that a correct procedure decide late: it must be
free to move already-allocated chores between bundles, and free to give up total
cost while doing so. Enforcing coverage \emph{structurally} rather than by price
meets that requirement --- if the copies of a chore simply cannot be placed on the
same agent, coverage holds by construction and no charge is needed --- and it is
the direction the constructions in progress take.
